\section{Erd\H{o}s Problem \#254}

\subsection*{1) ROUND-2 OBJECTIVE}
Path (C): isolate a single sharp obstruction whose resolution would settle the problem.

Round-1 extracted only very basic consequences from the hypotheses. The best Round-2 move is to prove a dyadic bootstrap theorem showing that, once one finite initial segment of $A$ creates a long enough interval of subset sums, the hypothesis
\[
|A\cap[1,2x]|-|A\cap[1,x]|\to\infty
\]
by itself propagates that interval forever. This would identify the exact missing bridge from the hypotheses to completeness.

\subsection*{2) Round-1 FOUNDATION USED}
I use only the following Round-1 facts.

\begin{itemize}
\item Writing
\[
B_j:=A\cap(2^j,2^{j+1}],
\qquad
m_j:=|B_j|,
\]
Round-1 showed that $m_j\to\infty$.
\item Round-1 also showed that the second hypothesis implies $\gcd(A)=1$, but the new theorem below does not actually need that fact.
\end{itemize}

\subsection*{3) NEW INSIGHT / TOOL (ROUND-2)}
The new tool is an \emph{interval absorption lemma} for subset sums.

It shows that the only thing still missing is a ``starter interval'' of subset sums inside one sufficiently large dyadic initial segment.

More precisely, I prove:

\begin{quote}
If for some large $J$ the finite set $A\cap[1,2^J]$ already has a subset-sum interval of length at least $2^{J+1}$, then the dyadic growth condition $m_j\to\infty$ forces $A$ to be complete.
\end{quote}

So the second hypothesis involving
\[
\sum_{n\in A}\{\theta n\}=\infty
\]
must, in any successful proof, be used only to manufacture such a starter interval.

\subsection*{4) ATTACK PLAN (ROUND-2)}
The exact Round-1 gap was the absence of any mechanism turning the two qualitative hypotheses into actual intervals inside the subset-sum set
\[
\Sigma(A):=\left\{\sum_{a\in F}a:\ F\subseteq A\text{ finite}\right\}.
\]

The new plan is:

\begin{itemize}
\item prove a general interval-absorption lemma for subset sums;
\item apply it block-by-block to the dyadic decomposition $A=\bigsqcup_j B_j$;
\item conclude that once one interval has length at least the next dyadic scale, all later dyadic blocks can be absorbed greedily.
\end{itemize}

This produces a single sharply formulated obstruction to completeness.

\subsection*{5) WORK (ROUND-2)}
For a finite set $F\subseteq \mathbf{N}$, write
\[
\Sigma(F):=\left\{\sum_{a\in E} a:\ E\subseteq F\right\}.
\]

\paragraph{Lemma 1 (interval absorption by one new element).}
Let $F\subseteq \mathbf{N}$ be finite, and suppose there is an interval
\[
[L,U]\subseteq \Sigma(F).
\]
Let $c\in \mathbf{N}$ with
\[
c\le U-L+1.
\]
Then
\[
[L,U+c]\subseteq \Sigma(F\cup\{c\}).
\]

\paragraph{Proof.}
The sums not using $c$ give the interval $[L,U]$, and the sums using $c$ give the translated interval
\[
[L+c,U+c].
\]
The assumption $c\le U-L+1$ is equivalent to
\[
L+c\le U+1,
\]
so these two intervals overlap or are adjacent. Hence their union is the full interval
\[
[L,U+c].
\]
Every sum in this union uses distinct elements because $c$ is a new element not in $F$.
\hfill$\square$

\medskip

\paragraph{Theorem 2 (dyadic bootstrap from one starter interval).}
Assume $m_j\to\infty$. Suppose there exists an integer $J$ such that:

\begin{itemize}
\item $m_j\ge 2$ for every $j\ge J+1$;
\item the finite set $A\cap[1,2^J]$ has a subset-sum interval
\[
[L_J,U_J]\subseteq \Sigma\bigl(A\cap[1,2^J]\bigr)
\]
of length
\[
W_J:=U_J-L_J+1\ge 2^{J+1}.
\]
\end{itemize}

Then every sufficiently large integer belongs to $\Sigma(A)$.

\paragraph{Proof.}
We induct on $j\ge J$.

Induction hypothesis: there exists an interval
\[
[L_J,U_j]\subseteq \Sigma\bigl(A\cap[1,2^j]\bigr)
\]
of length
\[
W_j:=U_j-L_J+1\ge 2^{j+1}.
\]
This is true for $j=J$ by assumption.

Now fix $j\ge J$ and write the next dyadic block as
\[
B_{j+1}=\{b_1<b_2<\cdots<b_t\},
\qquad t=m_{j+1}.
\]
Every element of $B_{j+1}$ satisfies
\[
2^j < b_i \le 2^{j+1} \le W_j.
\]
Therefore Lemma 1 applies to $b_1$, giving an interval of length $W_j+b_1$. Since $b_2\le 2^{j+1}\le W_j < W_j+b_1$, Lemma 1 applies again to $b_2$. Continuing in this way, we can absorb the entire block $B_{j+1}$ one element at a time, obtaining an interval
\[
[L_J,U_{j+1}]\subseteq \Sigma\bigl(A\cap[1,2^{j+1}]\bigr)
\]
of length
\[
W_{j+1}=W_j+\sum_{i=1}^{t} b_i.
\]
Because $t\ge 2$ and each $b_i>2^j$,
\[
\sum_{i=1}^{t} b_i > 2\cdot 2^j = 2^{j+1}.
\]
Hence
\[
W_{j+1}>W_j+2^{j+1}\ge 2^{j+1}+2^{j+1}=2^{j+2}.
\]
So the induction closes.

Thus for every $j\ge J$, the finite set $A\cap[1,2^j]$ has an interval of subset sums with the same left endpoint $L_J$ and with length at least $2^{j+1}$. As $j\to\infty$, these lengths tend to infinity, so
\[
[L_J,\infty)\subseteq \Sigma(A).
\]
Therefore every sufficiently large integer is a sum of distinct elements of $A$.
\hfill$\square$

\medskip

\paragraph{Corollary 3 (sharp obstruction to a counterexample).}
Assume the hypotheses of Problem \#254 and suppose $A$ is \emph{not} complete. Then for every sufficiently large $J$ the finite subset-sum set
\[
\Sigma\bigl(A\cap[1,2^J]\bigr)
\]
contains no interval of length $2^{J+1}$.

\paragraph{Proof.}
Round-1 gives $m_j\to\infty$, so there exists $J_0$ such that $m_j\ge 2$ for all $j\ge J_0$. If for some $J\ge J_0$ the set $\Sigma(A\cap[1,2^J])$ contained an interval of length $2^{J+1}$, then Theorem 2 would imply that $A$ is complete, contradiction.
\hfill$\square$

\medskip

\paragraph{Interpretation.}
The unresolved part of Problem \#254 is now concentrated in one sharply formulated finite statement:

\begin{quote}
Produce, from the two original hypotheses, one initial dyadic scale $2^J$ for which the finite set $A\cap[1,2^J]$ already has a subset-sum interval of length at least $2^{J+1}$.
\end{quote}

Once such a scale exists, the dyadic growth hypothesis alone finishes the proof.

\subsection*{6) ADVERSARIAL VERIFICATION}

\begin{itemize}
\item The absorption lemma depends on interval \emph{length}, not on the right endpoint alone. This was checked carefully: the overlap condition is exactly $c\le U-L+1$.
\item Distinctness of summands is preserved, because each absorbed element is added only once and is not already present in the finite set used at that step.
\item The theorem needs $m_j\ge 2$ eventually, not merely $m_j\ge 1$. This is why the hypothesis $m_j\to\infty$ is used at precisely that point.
\item The second hypothesis involving $\sum_{n\in A}\{\theta n\}$ plays no role in the bootstrap theorem itself; this is intentional, and it cleanly isolates where any future proof must use that hypothesis.
\end{itemize}

So the new theorem is internally consistent and pinpoints a single missing bridge rather than leaving the obstruction diffuse.

\subsection*{7) FINAL}
\textbf{UNRESOLVED (BUT STRICTLY ADVANCED).}

The strict Round-2 advance is a dyadic bootstrap theorem:

\begin{quote}
If one can obtain just one subset-sum interval of length at least the next dyadic scale inside some finite initial segment $A\cap[1,2^J]$, then the dyadic growth condition propagates this interval forever and forces completeness.
\end{quote}

Equivalently, any genuine counterexample would have to avoid such starter intervals at every sufficiently large dyadic scale. That is a much sharper obstruction than in Round-1.

\subsection*{8) COMPLETION ESTIMATE (MANDATORY)}
COMPLETION: 35\%

\subsection*{9) REFERENCES}
Only references actually used or checked.

\begin{itemize}
\item T. F. Bloom, \emph{Erd\H{o}s Problem \#254}, \texttt{https://www.erdosproblems.com/254}, accessed 2026-03-07.
\item J. W. S. Cassels, \emph{On the representation of integers as the sums of distinct summands taken from a fixed set}, Acta Sci. Math. (Szeged) 21 (1960), 111--124.
\end{itemize}
